\documentclass[11pt]{article}
\usepackage[margin=1in]{geometry}
\usepackage{amsmath,amssymb,amsthm}
\newtheorem{theorem}{Theorem}
\title{Adaptive support corrections and a sharp support obstruction}
\author{Joshua Mark}
\date{20 November 2026; research draft, priority unestablished}
\begin{document}\maketitle
Let $S=\{x_1,\ldots,x_s\}\subseteq\{\pm1\}^d$ consist of distinct
vectors. Put $k=1+d+\binom d2$ and let $A$ be the $k\times s$ matrix
whose rows are evaluations on $S$ of the characters of degrees zero,
one and two. Set $b=(1,0,\ldots,0)^T$. A probability vector $p$ is
unbiased and pairwise independent precisely when $Ap=b$.

\begin{theorem}[Actual positive correction with an explicit margin]
Suppose $A$ has full row rank, and define
\[
 R=A^T(AA^T)^{-1},\quad
 \gamma=\max_j\sum_a|R_{ja}|,\quad
 m=Ap-b,\quad \tau=\|m\|_\infty.
\]
For any probability vector $p$ with $p_{\min}\ge\gamma\tau$, the
vector $q=p-Rm$ is a nonnegative probability vector on the same
support and satisfies $Aq=b$. In particular this is an exact
distribution correction, not an unconstrained projection being
mistaken for a probability law. It obeys
\[
 \|q-p\|_\infty\le\gamma\tau,\quad
 d_{\rm TV}(p,q)\le s\gamma\tau/2,\quad
 C(q)\ge C(p)-2\gamma\tau.
\]
\end{theorem}
\begin{proof}
Full row rank gives $AR=I_k$, so $Aq=b$. The definition of $\gamma$
gives $|(Rm)_j|\le\gamma\tau$, proving positivity and both distance
bounds. The degree-zero equation of $Aq=b$ says $\sum q_j=1$.
Finally write $h=q-p$; then
$C(q)=C(p)+2\sum p_jh_j+\sum h_j^2\ge C(p)-2\|h\|_\infty$.
\end{proof}
All inverse entries, rank conditions and margins are exact rational
calculations for rational input weights. If $A$ is not full row rank,
one may instead solve on its column span, but first must verify
$b\in\operatorname{range}A$; small residual alone does not imply this.
The positivity margin must still be checked. The theorem gives no
license to omit either condition.

For $d=n-1$, $n\equiv1\pmod4$, $N=n+3$, the extra hypotheses
\[
 C(p)\ge1/N,\qquad \gamma\tau\le1/(2048N^4)
\]
place $q$ in the round6 extraction window and thus allow exact
Hadamard extraction. This is a sufficient criterion on a supplied
distribution and support, not a construction satisfying it in every
unknown order. The conditioning of the support enters explicitly
through $\gamma$, and tiny reservoir weights enter through $p_{\min}$.

\begin{theorem}[No adaptive reweighting on a minimal-size support]
Let $N\ge8$ be divisible by four, $d=N-4$, and $S$ have at most $N$
distinct sign vectors. If an unbiased pairwise-independent law exists
on $S$, its positive support has exactly $N$ members and its weights
are all $1/N$. Its augmented sign rows complete to a Hadamard matrix
of order $N$.
Consequently, if the uniform law on a given $N$-point set is not
isotropic, no reweighting on that same set can make it isotropic.
\end{theorem}
\begin{proof}
If the positive support has size $t\le N$, Cauchy--Schwarz gives
$C\ge1/t\ge1/N$. The residue-sensitive collision bound proved in
round5 gives $C\le1/N$. Hence $t=N$ and equality in Cauchy--Schwarz
forces uniform weights. The equality case of that same bound gives
the rank-three tetrahedral completion to a Hadamard matrix of order
$N$.
\end{proof}
This obstruction applies to every adaptive weight correction on the
specified support, rather than only the scalar Fourier buffers of
round7. Such a correction must change or enlarge a failed minimal
support. Conversely the full-row-rank correction theorem necessarily
uses $s\ge k$, which exceeds $N$ when $d=N-4\ge4$. Its hypotheses
therefore address genuinely richer supports and cannot covertly
circumvent the minimal-support obstruction.
\end{document}
